\chapter[Reinforcement Learning for Motor Control]{Reinforcement Learning\\for Motor Control}
\label{ch:rl}

Reinforcement learning optimizes behavior from interaction data under a specified
reward and transition process \cite{Sutton2018}. It does not remove the need to
define an objective. A model-free policy update need not differentiate an
explicit dynamics model, but its simulated experience still comes from a
particular mechanical model. A policy that exploits an inaccurate contact rule
or an arbitrary reward can score well without explaining a human golf swing.

\section{State, Observation, Action and Reward}

A Markov state contains enough information to specify the conditional next-state
distribution given the action. Joint positions alone generally omit velocity;
muscle-driven systems may also need activation and tendon states. An observation
can omit parts of the state, producing a partially observed control problem.
Neither the observation nor the action must equal generalized coordinates or
joint torque. State these mappings before selecting an algorithm.

For a discounted continuing objective,
$J(\pi)=\mathbb{E}_{\pi}[\sum_{k=0}^{\infty}\gamma^k r_k]$ with
$0\leq\gamma<1$, the reward scale, discount, sampling interval and action
constraints jointly define the problem. Changing the interval while holding
$\gamma$ and per-step rewards fixed generally changes the physical-time
objective. A continuous discount rate $\rho$ corresponds to
$\gamma=\exp(-\rho\Delta t)$; a running reward rate needs integration over
the step. These choices matter when comparing policies with different
simulation or controller rates.

Gymnasium's \texttt{reset} returns \texttt{(observation, info)};
\texttt{step} returns observation, reward, \texttt{terminated},
\texttt{truncated} and info. Spaces declare action/observation shapes, bounds
and dtypes, and an environment uses its local random generator
\cite{GymnasiumEnvironmentReview}. Interface compatibility does not imply
compatibility with every learning algorithm: for example, a continuous torque
space requires an algorithm or adaptation that supports continuous actions.

\section{An Executable Pendulum Benchmark}

The listing runs the maintained \texttt{Pendulum-v1} environment, tested here
with Gymnasium 1.3.0. Install \texttt{gymnasium==1.3.0}; rendering extras are
unnecessary. This is a reproducible rollout, not a trained policy or a claim
that the companion platform supplies a validated muscle-driven golf environment.

The benchmark uses a uniform rod with $m=L=1$, $g=10$, $\Delta t=0.05$,
$|\tau|\leq2$, and $\theta=0$ \emph{upright}. Its observation is
$(\cos\theta,\sin\theta,\omega)$, and before numerical clipping,
\begin{equation}
  \dot\omega=\frac{3g}{2L}\sin\theta+\frac{3\tau}{mL^2}.
\end{equation}
The factors follow from pivot inertia $mL^2/3$ and COM distance $L/2$.
Positive gravity acceleration near $\theta=0$ makes the upright equilibrium
unstable. This angle differs from the downward-zero convention in
Chapter~\ref{ch:param-estimation}; shifting the zero changes the gravity sign.

The implementation first updates angular velocity, clips it to $[-8,8]$,
then updates angle using that new velocity. This is a semi-implicit Euler
step with an additional velocity cap, not exact integration or an unconstrained
physical pendulum. The pre-step reward is
\begin{equation}
  r_k=-\left(\bar\theta_k^2+0.1\omega_k^2+0.001\tau_k^2\right),
  \qquad \bar\theta_k\in[-\pi,\pi).
\end{equation}
These benchmark weights define a scalar objective after choosing units;
squared torque is not mechanical work or muscle metabolic cost
\cite{GymnasiumPendulumReview,GymnasiumPendulumSourceReview}.

\begin{lstlisting}[language=Python, caption={A Seeded Pendulum Policy Rollout}]
from collections.abc import Callable
import gymnasium as gym
import numpy as np

def run_pendulum(
    policy: Callable[[np.ndarray], np.ndarray],
    seed: int, steps: int,
) -> dict:
    """Collect a continuing-task rollout with an external limit."""
    if (isinstance(seed, bool) or not isinstance(seed, int)
            or seed < 0):
        raise ValueError("seed must be a nonnegative integer")
    if (isinstance(steps, bool) or not isinstance(steps, int)
            or steps < 1):
        raise ValueError("steps must be a positive integer")
    env = gym.make("Pendulum-v1", max_episode_steps=steps)
    try:
        obs, _ = env.reset(seed=seed)
        observations = [obs.copy()]
        actions, rewards, terminated, truncated = [], [], [], []
        for _ in range(steps):
            action = np.asarray(policy(obs.copy()), dtype=np.float32)
            if not env.action_space.contains(action):
                raise ValueError("action must be finite (1,) in [-2,2]")
            obs, reward, terminal, cutoff, _ = env.step(action)
            observations.append(obs.copy())
            actions.append(action.copy())
            rewards.append(reward)
            terminated.append(terminal)
            truncated.append(cutoff)
            if terminal or cutoff:
                break
        return {
            "observations": np.asarray(observations),
            "actions": np.asarray(actions),
            "rewards": np.asarray(rewards),
            "terminated": np.asarray(terminated, dtype=bool),
            "truncated": np.asarray(truncated, dtype=bool),
        }
    finally:
        env.close()
\end{lstlisting}

For a zero-torque baseline, pass \texttt{lambda obs: np.array([0.0])}
as the policy, with seed 42 and 200 steps. The observation record includes the initial
state and one observation after each action; rewards correspond to the
pre-step state. The wrapper rejects actions outside the declared space instead
of relying on the benchmark's internal clipping. A stochastic policy needs its
own seeded generator; the environment seed does not seed an arbitrary policy.
Independent rollout seeds and uncertainty summaries are needed for a performance
comparison, not just one favorable return.

\section{Termination Is Not Truncation}

\texttt{Pendulum-v1} has no task-defined terminal state; its standard 200-step
limit is imposed by the \texttt{TimeLimit} wrapper. The listing sets that external
collection limit explicitly. Both flags stop collection, but an external
truncation does not say the continuing task has no future value. A one-step
target uses the final observation before reset:
\begin{equation}
  y_k=r_k+\gamma(1-\mathbf{1}_{\mathrm{terminated}})V(o_{k+1}).
\end{equation}
An intrinsically finite-horizon task is different: reaching its horizon is
termination, and the remaining time must be represented in a Markov state
\cite{GymnasiumTimeLimitsReview}. Algorithms and vectorized wrappers must also
distinguish the final observation from an automatically reset observation.

\begin{lstlisting}[language=Python, caption={A Continuing-Task Bootstrap Target}]
def bootstrap_target(
    reward: float, discount: float,
    next_value: float, terminated: bool,
) -> float:
    """Use terminated, not terminated-or-truncated, as the mask."""
    values = np.asarray([reward, discount, next_value], dtype=float)
    if values.shape != (3,) or not np.isfinite(values).all():
        raise ValueError("reward, discount and value must be finite scalars")
    if not 0 <= discount <= 1:
        raise ValueError("discount must be in [0,1]")
    if not isinstance(terminated, (bool, np.bool_)):
        raise ValueError("terminated must be boolean")
    return float(reward + discount * (0.0 if terminated else next_value))
\end{lstlisting}

For $r=2$, $\gamma=0.9$ and $V(o')=5$, an external cutoff gives $y=6.5$;
a true terminal transition gives $y=2$. This distinction can bias learning even
when the simulator and reward are otherwise correct. The helper permits
$\gamma=1$ for finite-horizon or other explicitly justified formulations;
it does not assert convergence of an undiscounted infinite sum.

\section{From a Benchmark to a Golf Hypothesis}

PPO is one policy-optimization method, not an assurance of global optimality,
constraint satisfaction or transfer to a golfer \cite{schulman2017}. Learning
curves should report environment transitions, evaluation seeds and uncertainty,
as well as reward and policy settings. Comparisons with trajectory optimization
need matched dynamics, initial conditions, horizons, constraints and objectives;
Chapter~5 supplies a local quadratic core, not an end-to-end iLQR baseline.

In golf, rewarding speed at a loosely defined impact can favor missing the
ball, an incorrect face orientation or infeasible contact. Specify the strike
event, club and ball geometry, launch-related quantities, feasible actuator
states, joint/contact constraints and failure conditions. Penalty rewards do
not by themselves enforce hard constraints. Domain randomization explores the
chosen distribution of model variation; it cannot certify robustness to omitted
mechanics or identify the human nervous system's objective.

An interpretable comparison therefore links Chapters~\ref{ch:param-estimation}
and~\ref{ch:visualization}: identify what the data constrain, declare the
remaining model uncertainty, optimize an explicit task, then test motion,
forces and energy on held-out conditions. Similar club paths can arise from
different policies and actuator assumptions. A successful learned swing is a
constructive possibility within the model, not proof that a golfer uses that
policy.

\section*{Exercises}
\begin{enumerate}
  \item Reproduce one benchmark transition and its pre-step reward from the
  stated equations. Explain how the velocity cap changes the mechanical model.
  \item Compare true termination and external truncation for the numerical
  target above. State which observation a vectorized learner should bootstrap.
  \item Specify a golf reward and a separate feasibility test. Give one
  physically undesirable behavior that the reward alone could admit.
\end{enumerate}
